\documentclass[../../main.tex]{subfiles}
\begin{document}
\begin{oframed}
  \subfile{problem}
\end{oframed}

{\bf [解答]} $a\not=0$である．まず，直線の式で$x,y>0$であるから，$0<x<4$となる．方程式から$y$を消去した
\begin{align}
  4-x=-ax^2+bx \Longleftrightarrow ax^2-(b+1)x+4=0
\end{align}
が$0<x<4$に重解を持つ．判別式$D$として
\begin{align}
  \begin{cases}
    D=0 \\
    0<\cfrac{b+1}{2a}<4
  \end{cases}
  \Longleftrightarrow
  \begin{cases}
    (b+1)^2=16a \\
    0<\cfrac{b+1}{2a}<4
  \end{cases}
\end{align}
である．
$a\neq 0$に注意して第一式から$a>0$であるから，
\begin{align}
  &
  \begin{cases}
    (b+1)^2=16a \\
    0<b+1<8a
  \end{cases} \\
  \Longleftrightarrow
  &
  \begin{cases}
    (b+1)^2=16a \\
    0<b+1<\cfrac{(b+1)^2}{2}
  \end{cases} \\
  \Longleftrightarrow
  \begin{cases}
    (b+1)^2=16a \\
    1<b
  \end{cases} \label{eq:1}
\end{align}
この時，放物線の概形は\cref{fig:parabola_and_line}のようになる.

\begin{figure}[htb]
  \centering
  \begin{tikzpicture}[x=1.35cm, y=0.95cm, >=stealth, every node/.style={font=\small}]
    % 領域 S (斜線パターン)
    \fill[pattern=north east lines, pattern color=blue!35]
    (0, 0) -- plot[domain=0:3, samples=50] (\x, {-\x*\x + 3*\x}) -- cycle;

    % 座標軸
    \draw[->, thick, gray!80!black] (-0.6, 0) -- (4.8, 0) node[right] {$x$};
    \draw[->, thick, gray!80!black] (0, -0.6) -- (0, 4.8) node[above] {$y$};
    \node[below left=2pt] at (0, 0) {$O$};

    % 直線 x + y = 4
    \draw[line width=1.1pt, orange!85!black, domain=-0.2:4.3]
    plot (\x, {4 - \x});
    \node[orange!90!black, above right=1pt, font=\footnotesize] at (3.2, 1.0) {$x+y=4$};

    % 直線の切片
    \node[below=2pt] at (4, 0) {$4$};
    \node[left=2pt] at (0, 4) {$4$};

    % 放物線 y = -ax^2 + bx (a=1, b=3 の例)
    \draw[line width=1.3pt, blue!85!black, domain=-0.25:3.25, samples=60]
    plot (\x, {-\x*\x + 3*\x});
    \node[blue!85!black, above left=1pt, font=\footnotesize] at (1.2, 2.3) {$y=-ax^2+bx$};

    % 放物線の x 切片
    \node[below left=2pt] at (3, 0) {$\frac{b}{a}$};

    % 接点 P (2,2)
    \coordinate (P) at (2, 2);
    \draw[densely dotted, thick, gray!70!black] (2, 0) node[below=2pt] {$2$} -- (P) -- (0, 2) node[left=2pt] {$2$};
    \node[circle, fill=red!85!black, inner sep=1.3pt] at (P) {};
    \node[above right=2pt, red!85!black, font=\footnotesize] at (P) {接点 $(2,2)$};

    % 領域ラベル S
    \node[blue!85!black, font=\large] at (1.0, 0.8) {$S$};
  \end{tikzpicture}
  \caption{放物線 $y=-ax^2+bx$ と直線 $x+y=4$ が第一象限で接する様子（$a=1, b=3$ の場合）と囲まれる領域 $S$}
  \label{fig:parabola_and_line}
\end{figure}

題意の面積$S$として
\begin{align}
  S&=\int_0^{b/a}(-ax^2+bx)dx \\
  &=\dfrac{a}{6}\left(\dfrac{b}{a}\right)^3=\dfrac{b^3}{6a^2}
\end{align}
である．以下\cref{eq:1}の下でこの最大値を考えれば良い．\cref{eq:1}を代入して
\begin{align}
  S&=\dfrac{b^3}{6}\dfrac{16^2}{(b+1)^4} \nonumber\\
  &=\dfrac{128}{3}\dfrac{b^3}{(b+1)^4}\label{eq:2}　
\end{align}
ここで$b>1$の時，AM-GMから
\begin{align}
  \dfrac{3b^3}{(b+1)^4}&\le\left(\dfrac{1}{4}\left(\dfrac{3}{b+1}+3\dfrac{b}{b+1}\right)\right)^4 \\
  &=\left(\dfrac{3}{4}\right)^4 \\
  \therefore &\dfrac{b^3}{(b+1)^4}\le \dfrac{3^3}{4^4}
\end{align}
である．等号成立は$\cfrac{3}{b+1}=\cfrac{b}{b+1}\Longleftrightarrow b=3$の時．これは\cref{eq:1}を満たす．従って\cref{eq:2}
に代入して
\begin{align}
  &S\le\dfrac{128}{3}\dfrac{3^3}{4^4}=\dfrac{9}{2}
  &\therefore \max S=\dfrac{9}{2}
\end{align}
である．等号成立は\cref{eq:1}から$(a,b)=(1,3)$の時である．$\cdots$(答)

\newpage
\end{document}
